% The spoken body of the homily for the Twenty-first Sunday after Pentecost.
% Continuous preaching, ready to read aloud: no heading, no direction to the
% preacher, no citation inside the speech. The loci, the relation to the
% reviewed readings, the audience, the word count and the pace estimate stand
% in the terminal note, which no one reads out. The blank lines are the
% paragraph breaks a speaker recovers his place by, and the vertical spaces
% divide the speech into its movements: the same plea said twice; the two
% debts; where the debt was forgiven, the Cross and the empty tomb, received
% in Baptism with Augustine; the plea heard and refused; the debt that
% returns, with Chrysostom and Aquinas's distinction; the wrong throat, with
% the Epistle, Chrysostom on Ephesians, the Introit and the Offertory; "from
% your hearts", with Augustine's tongue and heart and his
% surgeon; the Communion's plea, with Augustine's patience of God; the altar,
% with the Secret, Augustine's beggars and the Postcommunion; the week's
% response, with the Lord's Prayer and the Collect; and the return to the
% courtyard. The honorific is tied to the name so that "St" never ends a line.
\raggedcolumns

``Have patience with me, and I will pay thee all.''

Twice in today's Gospel a man falls on his knees and says those words. The
first time, he says them to a king, and he owes the king ten thousand talents.
The second time, they are said to him, by a man who owes him a hundred pence.
The words are the same, word for word. Only the man who hears them has changed
places.

The first time, those words save him. The second time, he refuses them.

So the Gospel asks each of us: what does a forgiven man owe the man who has
wronged him?

\medskip

The whole parable hangs on the distance between the two debts. A hundred
pence is not nothing. It was about a hundred days' wages for a
labourer. The Gospel does not tell you that a wrong done to you does not
matter.

But ten thousand talents is another kind of number. It is a debt no servant
could ever pay. And the servant asks only for time. He is given the whole
debt.

Who is that servant? St~Augustine answers that every one of us is both men in
this story at once. We owe God a debt we cannot pay. And somewhere there is
someone who owes us.

\medskip

In the parable the king forgives with a word. We know what that word cost.

Our debt was not cancelled by God looking the other way. The King's own Son
took it on himself. He paid it on the Cross, in his own blood, and the Father
raised him from the dead, so that the pardon stands for ever. That is where
the ten thousand talents were forgiven: not in a story, but on Calvary and at
the empty tomb.

And that pardon reached each of us at the font. St~Augustine says it plainly:
``every man begins from Baptism; he goes out free, the `ten thousand talents'
are forgiven him; and when he goes out, he will soon find some fellow-servant
his debtor.''

\medskip

Now watch the forgiven man go out. He goes ``straightway'', Chrysostom notices, with the gift still fresh upon
him. He finds his fellow servant, takes him by the throat, and says, ``Pay what
thou owest.'' And his fellow servant falls down and says the very words: ``Have
patience with me, and I will pay thee all.''

This is the moment on which the whole Gospel turns. The servant is hearing his
own prayer. The words that had just saved him are now in another man's mouth.
Chrysostom says that he ``did not recognize so much as the harbor by which he
escaped shipwreck.''

That is his sin. Not that he wanted his money back, but that he forgot what
had been done for him. A man who remembered the ten thousand talents could not
have kept his hand on that throat.

\medskip

Then the king calls him back. ``Thou wicked servant, I forgave thee all the
debt, because thou besoughtest me.'' And he hands him over ``until he paid all
the debt.''

All the debt. The debt that was forgiven comes back. Who took the pardon away?
Chrysostom answers, ``it was not he that revoked it, but the one who had
received it.''

St~Thomas Aquinas helps us hear this rightly. A grudge does not
undo your baptism. The sins God has forgiven do not come back as the same sins.
But ingratitude makes a man a debtor again. To receive mercy and refuse it to
another is to put yourself back under the whole weight of the debt. ``So also
shall my heavenly Father do to you, if you forgive not every one his brother
from your hearts.''

\medskip

Why does the forgiven servant do it? The Epistle gives part of the answer.
St~Paul writes, ``our wrestling is not against flesh and blood; but against
principalities and powers, against the rulers of the world of this darkness.''

Explaining this passage, St~John Chrysostom says something startling. ``If any
one hath an enemy, if any one hath been wronged by him, if any one is
exasperated, let him collect together all that wrath, all that fierceness, and
pour it out upon the head of the devil.''

The servant had his hands on the wrong throat. The man who owed him a hundred
pence was not his enemy. He was a fellow servant, a fellow debtor, standing
before the same king. The enemy was the one who wanted the grudge.

Chrysostom does not call the wrong done to you nothing. He tells you where to
aim your anger. And this Mass
shows how far the real enemy can go. It began by confessing, ``All things,
O Lord, are in thy power: and no one can resist thy will.'' Its Offertory tells
how Satan asked to tempt Job, and power was given him by the Lord. Given, and
measured.

\medskip

There is one more phrase in the Gospel's last sentence: ``from your hearts.''

St~Augustine will not let us forgive with the mouth alone. ``Say not with the
tongue, `I forgive,' and put off to forgive in the heart.''

Forgiveness is not pretending. A boy who keeps offending, Augustine says, may
still be corrected. ``For what is so kind and gentle as the surgeon with his
knife?'' What must go is the account we keep against our brother in the heart.
Forgiveness from the heart forbids rancour. It does not forbid correction.

\medskip

Later in this Mass, the Communion verse will ask a hard question: ``when wilt
thou judge those who persecute me?'' Has the Church forgotten the Gospel by
then?

St~Augustine, preaching on the Lord's Prayer, answers it. ``A man lusteth to
be avenged, and Christ is not yet avenged, the holy martyrs are not yet
avenged. Still doth the patience of God wait, that the enemies of Christ, the
enemies of the martyrs, may be converted.'' And so, he says, ``Forgive
therefore; from the heart forgive.''

So we may hand every wrong over to God's judgment; it is safe there. And his
patience waits for the conversion of the one who wronged us, as it once waited
for ours.

\medskip

That patience has a place where we can see it. It is this altar. Over the
bread and wine the Church will pray: ``Mercifully receive, O Lord, these offerings, by which thou art pleased
to be appeased, and restore us to salvation by thy powerful goodness.'' What is
offered here is not a second payment beside the Cross. It is the one sacrifice
of the Cross, made present among us. The price of our pardon is laid on this
altar.

Then come to receive as the servant came to the king: on your knees, with
nothing in your hands to pay. St~Augustine says that when we pray we are
all God's beggars, and that what we beg of God is Christ, the living Bread. At
this altar the beggar is given what he asks: not a delay, but the Lord himself,
crucified and risen. The Church calls him ``the food of immortality'', and asks
that the pure mind may follow after what the mouth has received. To receive
him is to share his death and his rising, and to follow him is to go
out as he came out to us.

\medskip

What, then, shall we do this week? First, name the debt. Most of us know at once whose throat we are holding. Name
that person before God.

Second, before Communion the prayer Christ taught us is said at this altar.
Mean its condition: ``Forgive us our debts, as we forgive our debtors.''

Third, when the anger rises, aim it where Chrysostom aims it: not at your
brother, but at the enemy who wants you to keep the grudge.

The Collect prayed, ``Preserve thy family, O Lord''. We are that family, a
household of servants, every one of us forgiven at the same price. Let no one who has received the Lord at this altar
go out and take his brother by the throat.

\medskip

Come back, then, to the king's house. A man walks out free. Before him a fellow
servant kneels and says, ``Have patience with me, and I will pay thee all.''

This time, let him know the words. They are the words by which he was saved.
Let him take his hand from that throat, and lift his brother up. The King has
forgiven him ten thousand talents at the price of his Son's blood, and feeds him
at his own table. A hundred pence is not too much to forgive, for a man whom
the King has made so rich.
